\chapter{Real-замыкание проходов и граница рамированного вложения}
\label{ch:v6-single-thread-framed-winding-embedding-gate}

\section{Необходимое условие замкнутой пространственной нити}

Целочисленный подъём предыдущей главы построил один абстрактный цикл, но
ещё не проверил его пространственное замыкание. Для любой замкнутой
кусочно-гладкой кривой $X(s)$ выполнено
\begin{equation}
 \oint dX=\oint \dot X(s)\,ds=0.
 \label{eq:v6-single-thread-zero-total-tangent}
\end{equation}
Поэтому ориентированная сумма всех проходов должна обращаться в нуль.

Если все $K_a$ проходов направлены только вдоль $+n_a$, рабочие неравные
кратности не удовлетворяют этому условию. Это тот же ненулевой полярный
ток, который был обнаружен первым моментным гейтом.

\section{Real-удвоение замыкает ориентированную сумму}

Проект уже содержит сопряжённые ориентации $E/E^*$ и линии $L/L^*$.
Возьмём для каждой оси одинаковое число прямых и обратных проходов:
\begin{equation}
 K_a^+=K_a^-=K_a.
\end{equation}
Для лучшего целочисленного приближения это даёт
\begin{equation}
 K^+=K^-=(22696855,829548,829548,829548)
\end{equation}
и общее число ориентированных посещений
\begin{equation}
 M_{\rm Real}=50370998.
\end{equation}
Тогда
\begin{equation}
 \sum_a(K_a^+-K_a^-)n_a=0
\end{equation}
точно, а нормированный физический первый момент имеет численную норму
$4.21\cdot10^{-17}$.

Второй момент нечувствителен к обращению $n\mapsto-n$, поэтому поле $Q_*$
сохраняется с нулевым остатком.

\section{Как сохранить нечётный тетраэдрический момент}

Обычный третий момент встречных Real-проходов сокращается. То же верно для
канонического двойственного момента без дополнительной ориентационной
метки. Но на ориентированных стрелках уже существует знак степени:
\begin{equation}
 \chi(E)=+1,
 \qquad
 \chi(E^*)=-1.
\end{equation}
Пусть физическая касательная равна $t_{a,s}=s n_a$, где $s=\pm1$, а
двойственная касательная --- $d_{a,s}=R^{-1}t_{a,s}$. Рассмотрим
\begin{equation}
 \mathcal T_{\rm Real}
 =\operatorname{STF}\sum_{a,s}\frac{w_a}{2}\,s\,
 \operatorname{Sym}(d_{a,s}\otimes t_{a,s}\otimes t_{a,s}).
\end{equation}
Поскольку и смешанный тензор, и Real-характер меняют знак при
$s\mapsto-s$, их произведение чётно.

\begin{theorem}[Совместимость замыкания, $Q$ и $T$]
Равновесное Real-удвоение обращает физический первый момент в нуль,
сохраняет второй момент $Q_*$ и при ориентационном характере $\chi=s$
воспроизводит
\begin{equation}
 \mathcal T_{\rm Real}=\frac34T_*.
\end{equation}
Численный остаток равен $4.49\cdot10^{-16}$.
\end{theorem}

Это сильная кинематическая совместимость. Однако проектный
ориентационный функтор пока назначает $L/L^*$ операторным стрелкам, а не
доказанному физическому наблюдаемому фундаментальной нити. Поэтому знак
$\chi=s$ остаётся каноническим кандидатом, но не завершённым родительским
выводом.

\section{Замкнутый маршрут ещё не является вложенной трубкой}

Из сбалансированного набора шагов легко построить замкнутое слово: пройти
произвольную последовательность прямых шагов, а затем обратную
последовательность противоположных шагов. Но такая конструкция в точности
повторяет уже пройденные рёбра. Для нити ненулевой толщины это означает
наложение трубки самой на себя.

Следовательно, выполнены разные уровни утверждения:
\begin{enumerate}
 \item абстрактный цикл без ветвления: пройден;
 \item нулевая суммарная касательная: пройдена;
 \item простой самонепересекающийся цикл с тем же словом: не построен;
 \item положительный радиус трубки: не выведен;
 \item динамический барьер пересоединения: не вычислен.
\end{enumerate}

Стандартная геометрия рамированной замкнутой кривой различает
самозацепление, скручивание рамки и изгиб осевой линии; их связь выражает
формула Кэлугэряну--Уайта
\cite{single-thread-white1969,single-thread-dennis-hannay2005}. Но в
проекте ещё не построены нормальная рамка вдоль всего winding-слова и её
числа $Lk$, $Tw$, $Wr$.

\section{Почему известная ширина вихря не закрывает вопрос}

Эффективная вихревая нить Тома VI имеет конечную энергетическую ширину
$0.57328995$. Однако пересмотр приоритета уже отделил эту полевую трубку
от фундаментальной нити. Перенос её ширины в микроскопическую модель был
бы новым словарём без родительского масштаба.

Запрет самопересечения требует либо явного условия инъективности
$X(s)\ne X(s')$, либо положительной энергии исключённого объёма, как в
континуальных моделях самонепересекающихся цепей
\cite{single-thread-edwards1965}. Ни один такой член текущим спектральным
родителем не выведен.

\begin{nogocriterion}
Нельзя считать хопфов класс $c_1=1$ автоматически рамкой всей
пространственной нити. Он задаёт топологическую линию над базой, но не
разносит $5\cdot10^7$ соседних витков на положительное расстояние и не
создаёт энергетический барьер пересоединения.
\end{nogocriterion}

\section{Вердикт и следующий тест}

Real-удвоение снимает важное возможное противоречие: замкнутая нить может
иметь нулевой физический поток, сохраняя $Q_*$ и ориентационно-взвешенный
$T_*$. Поэтому гипотеза одной нити не опровергнута.

Но геометрическая неразрывность ещё не выведена. Следующий гейт должен
проверить, возникает ли из уже существующих кривизны, рангового барьера
или конечной ширины положительный excluded-volume функционал, расходящийся
при самопересечении и тем самым запрещающий пересоединение.

Нулевая гипотеза: родительская кривизна или барьер потери ранга создаёт
неотрицательную энергию взаимного сближения различных winding-участков без
нового коэффициента.

Стоп-критерий: если расстояние между несоседними участками не входит ни в
один существующий оператор, запрет пересоединения является новой аксиомой.

\begin{protocol}
\begin{enumerate}
 \item нулевая ориентированная сумма проходов: \textbf{да после Real-удвоения};
 \item общее число ориентированных посещений: \textbf{$50370998$};
 \item физический первый момент: \textbf{ноль};
 \item второй момент $Q_*$: \textbf{сохранён};
 \item нечётный момент без Real-характера: \textbf{сокращается};
 \item момент с $\chi=s$: \textbf{$3T_*/4$};
 \item родительский вывод $\chi=s$: \textbf{не завершён};
 \item замкнутый абстрактный маршрут: \textbf{существует};
 \item самонепересекающееся вложение: \textbf{не построено};
 \item фундаментальный радиус трубки: \textbf{не выведен};
 \item барьер пересоединения: \textbf{не выведен};
 \item следующая проверка: \textbf{исключённый объём и пересоединение}.
\end{enumerate}
\end{protocol}

Машинный сертификат сохранён в
\path{s2t_v6_single_thread_framed_winding_embedding_gate_results.json}.

\begin{thebibliography}{9}
\bibitem{single-thread-white1969}
J.~H.~White,
\emph{Self-Linking and the Gauss Integral in Higher Dimensions},
American Journal of Mathematics 91 (1969), 693--728.

\bibitem{single-thread-dennis-hannay2005}
M.~R.~Dennis, J.~H.~Hannay,
\emph{Geometry of C\u{a}lug\u{a}reanu's theorem},
Proceedings of the Royal Society A 461 (2005), 3245--3254;
arXiv:math-ph/0503012.

\bibitem{single-thread-edwards1965}
S.~F.~Edwards,
\emph{The statistical mechanics of polymers with excluded volume},
Proceedings of the Physical Society 85 (1965), 613--624.
\end{thebibliography}